\subsection{A CM descent makes the harmonic cover a finite-support target}

\begin{statement}{A.160}{Proved}
{Full Tate--Shafarevich finiteness and the exact cuspidal range of the cover}
Retain the constants, curves and actual maps of A.150 and A.154:
\[
 \begin{gathered}
 K=\Q(t),\quad t^2=2,\quad \epsilon=1+t,\quad
 M=325596479663,\quad b=1146\epsilon^2M^2,\\
 E_b/K:\ y^2=x^3+bx,\qquad
 \overline D:\ aW^4=T(T-1)(T-\tfrac12),\quad
 a=573\left(\frac{M}{2\cdot191^2}\right)^2.
 \end{gathered}
\tag{A.160.1}
\]
The full group \(\Sha(E_b/K)\) is finite.  Its two-primary
part remains the exactly computed group \((\Z/2)^2\) of
A.154; no order of its odd part is asserted.  Put
\(A=\operatorname{Res}_{K/\Q}E_b\).  There are exact
identifications
\[
 \begin{gathered}
 A(\Q)=E_b(K)=\{O,T_0\},\\
 \Sha(A/\Q)=\Sha(E_b/K),\\
 T_\ell\Sha(A/\Q)=0\quad\text{for every prime }\ell.
 \end{gathered}
\tag{A.160.2}
\]
The adjoint of the positive Prym map is a nonconstant
\(\Q\)-morphism \(g:\overline D\to A\), and
\[
 g^{-1}(O_A)=\{P_0,P_\infty\},\qquad
 g^{-1}(T_A)=\{P_1,P_{1/2}\}.
\tag{A.160.3}
\]
These are equalities of reduced finite schemes.

Every full Selmer section of a hyperbolic open subcurve
\(X\subseteq\overline D\) over \(\Q\) is the point section
of one of the four points in (A.160.3) belonging to \(X\),
or belongs to the cusp packet at one of them outside \(X\).
The identical classification holds over \(K=\Q(\sqrt2)\)
for every hyperbolic open subcurve of \(\overline D_K\),
with the same four rational boundary points.  Over \(K\)
the actual target is \(E_b\) and the morphism is \(q_+\).
In particular every full Selmer section of
\(D^\circ=\overline D\setminus\{P_0,P_1,P_{1/2},P_\infty\}\)
is cuspidal.  Every full Selmer section of the tripod
whose combined class satisfies
\[
                   2\kappa_8(t_U)+\kappa_8(1-t_U)=[64a^2]
\tag{A.160.4}
\]
is cuspidal.  Here \(t_U=4T(1-T)\); the use of \(t_U\)
distinguishes the tripod coordinate from the fixed element
\(t=\sqrt2\) in (A.160.1).

There is also an all-coefficient quotient statement.  With
\(H\) and \(\mathcal I\) as in A.157, set
\[
 N_{\rm all}=\ker\{H\longrightarrow(H/\mathcal I)^{\rm ab}\}.
\tag{A.160.5}
\]
This subgroup is arithmetic normal, and
\[
 \Pi/N_{\rm all}\simeq
 \bigl(\widehat\Z[V]/(V^6+V^4+V^2+1)\bigr)\rtimes C_8,
       \qquad x\mapsto(1,v^2),\quad y\mapsto(0,v).
\tag{A.160.6}
\]
A continuous section of \(E_U/N_{\rm all}\) in the class
(A.160.4) which locally matches actual point or cusp
sections has no ordinary localization at any finite prime.
The same assertion follows for the full metabelian quotient
through its actual map to (A.160.6).  The identification of
one global cusp packet above is asserted for full sections;
the quotient assertion specifies the local images just stated.
\end{statement}

\begin{proof}
\emph{A literal model over its imaginary quadratic CM field.}
Put
\[
 F=K(i),\qquad K_0=\Q(i),\qquad
 \delta=2(1+i),\qquad k=\frac{1+i+t}{2(1+i)},\qquad
 c=9168iM^2.
\tag{A.160.7}
\]
The identities
\[
 (1+i+t)^2=\delta\epsilon,\qquad
 \epsilon=\delta k^2,\qquad \delta^2=8i,\qquad b=ck^4
\tag{A.160.8}
\]
give an actual \(F\)-isomorphism
\[
 C_F\xrightarrow{\sim}(E_b)_F,\qquad
 C/K_0:\ y^2=x^3+cx,\qquad (x,y)\longmapsto(k^2x,k^3y).
\tag{A.160.9}
\]
Its inverse divides the coordinates by \(k^2,k^3\).
Both \(C\) and its quadratic two-twist
\(C^{(2)}:y^2=x^3+4cx\) are defined over \(K_0\).
The automorphism \((x,y)\mapsto(-x,iy)\), defined over
\(K_0\), has square \([-1]\).  It embeds \(\Z[i]\) in
their endomorphism rings.  In characteristic zero the
endomorphism algebra of an elliptic curve containing this
imaginary quadratic field is that field; the maximality of
the order \(\Z[i]\) identifies the endomorphism ring with
\(\Z[i]\).  Thus these are elliptic curves over their CM
field with all the required CM endomorphisms defined there.

\emph{The two restriction--corestriction comparisons.}
We give the actual quadratic maps used twice below.
For the displayed equation \(B:y^2=x^3+dx\) over a field
\(k_0\), and a quadratic extension \(L=k_0(\sqrt e)\),
write \(B^{(e)}:y^2=x^3+e^2dx\), and let \(\sigma\)
be the nontrivial automorphism of \(L/k_0\).  The map
\[
 j_e:B^{(e)}_L\xrightarrow{\sim}B_L,
       \qquad (x,y)\longmapsto
                  \left(\frac{x}{e},\frac{y}{e\sqrt e}\right)
       \quad\text{satisfies }\sigma(j_e)=-j_e.
\tag{A.160.10}
\]
In the split \(L\)-coordinates of the Weil restriction,
the following maps are defined over \(k_0\):
\[
 \begin{aligned}
 \Phi:\operatorname{Res}_{L/k_0}B_L&\longrightarrow
                                  B\times B^{(e)},
     &(P,Q)&\longmapsto(P+Q,j_e^{-1}(P-Q)),\\
 \Psi:B\times B^{(e)}&\longrightarrow
                                  \operatorname{Res}_{L/k_0}B_L,
     &(R,S)&\longmapsto(R+j_eS,R-j_eS).
 \end{aligned}
\tag{A.160.11}
\]
Indeed \(\sigma\) exchanges \(P,Q\) and changes the sign
of \(j_e\), so the formulas are equivariant.  Both
compositions are multiplication by two.  On \(k_0\)-points
the first map reads
\(P\mapsto(P+\sigma P,j_e^{-1}(P-\sigma P))\).
On Tate modules after tensoring with \(\Q_2\), the inverse
of \(\Phi\) is \(\frac12\Psi\); the two resulting
projectors are \((1+\sigma)/2\) and \((1-\sigma)/2\),
with the latter identified with the twist by \(j_e\).
These same maps act on the Kummer sequences and their local
restrictions.  In particular they preserve the Selmer local
conditions, so the projector statement also holds on their
rational two-adic Selmer spaces.

For a class \(\xi\in\Sha(B/L)\), ordinary Galois
corestriction and the same twist give the more precise
integral identities
\[
 \begin{aligned}
 r_+(\xi)&=\operatorname{cor}_{L/k_0}\xi
                                      \in\Sha(B/k_0),\\
 r_-(\xi)&=\operatorname{cor}_{L/k_0}(j_e^{-1})_*\xi
                                      \in\Sha(B^{(e)}/k_0),\\
 \operatorname{res}r_+(\xi)&=(1+\sigma)\xi,\qquad
 (j_e)_*\operatorname{res}r_-(\xi)=(1-\sigma)\xi.
 \end{aligned}
\tag{A.160.12}
\]
The last sign follows from \(\sigma(j_e^{-1})=-j_e^{-1}\).
Corestriction is the sum over the two cosets, giving these
equalities at the level of cocycles modulo coboundaries.
Its local restriction is the sum of the local
corestrictions over places above a given place.  A locally
trivial class therefore remains locally trivial; this
proves the stated maps on \(\Sha\), including at the
archimedean places.

First apply these formulas with \(k_0=K,L=F,e=-1\).
The twist \(E_b^{(-1)}\) has literally the same equation
as \(E_b\), and \(j_{-1}(x,y)=(-x,iy)\).
A.154 says that both groups on the right of (A.160.12)
have two-primary part killed by two.  Hence
\[
             4\Sha(E_b/F)[2^\infty]=0.
\tag{A.160.13}
\]
For points, both entries of \(\Phi(P)\) belong to
\(E_b(K)=\{O,T_0\}\); (A.160.11) similarly gives
\([4]E_b(F)=0\).  Thus \(E_b(F)\) is finite, being a
subgroup of the finite group \(E_b[4](\overline F)\).

For every number field and fixed positive integer \(n\),
the multiplication Selmer group is finite.  Here the usual
finiteness argument applies directly: away from the primes
over \(n\) and of bad reduction, a finite Kummer class is
unramified.  After adjoining the finite \(n\)-torsion,
the resulting cocycles cut out extensions of bounded degree
unramified outside that fixed finite set.  There are
finitely many such extensions: local extensions of bounded
degree at the omitted primes have bounded discriminant,
and Hermite--Minkowski applies.  Each finite extension has
only finitely many cocycles with values in the fixed finite
module.  The finite multiplication sequence then makes
\(\Sha[n]\) finite as well.  Consequently (A.160.13)
proves that \(\Sha(E_b/F)[2^\infty]\) is finite.
The discrete \(2^\infty\)-Selmer group is finite too,
since its point contribution \(E_b(F)\otimes\Q_2/\Z_2\)
is zero.

Via (A.160.9), the same conclusions hold for \(C/F\).
The map \((x,y)\mapsto(x/2,y/(2t))\) identifies
\(C^{(2)}_F\) with \(C_F\).  Restriction from \(K_0\)
has kernel contained in
\[
 H^1(\operatorname{Gal}(F/K_0),C(F))
 \quad\text{or}\quad
 H^1(\operatorname{Gal}(F/K_0),C^{(2)}(F)),
\tag{A.160.14}
\]
respectively, by the cocycle inflation--restriction
sequence.  These are finite groups because their coefficient
groups are finite.  The point groups over \(K_0\) are
subgroups of the finite point groups over \(F\).
It follows that both
\(\operatorname{Sel}_{2^\infty}(C/K_0)\) and
\(\operatorname{Sel}_{2^\infty}(C^{(2)}/K_0)\) are finite.
In particular their \(\Z_2\)-coranks are zero.  This
derives the needed integral finiteness before using the
rational projectors of (A.160.11).

\emph{The exact primary theorems and all primes.}
Theorem 1.1 of the opened author version of
Burungale--Tian applies to an elliptic curve defined over
an imaginary quadratic field and having CM by an order
of that field, at every prime \(p\).  It sends vanishing
\(p^\infty\)-Selmer corank to nonvanishing of the central
complex \(L\)-value
\cite[Theorem 1.1, p.~1]{LegBurungaleTian2026}.
Apply it with \(K_0=\Q(i)\) and \(p=2\) to the two
curves just proved to satisfy its hypotheses.  This gives
\[
                  L(C/K_0,1)\ne0,
                  \qquad L(C^{(2)}/K_0,1)\ne0.
\tag{A.160.15}
\]
The first assertion of Rubin's Theorem A, opened in the
original printed page 527, gives finiteness of the entire
\(\Sha\) for each of these curves
\cite[Theorem A, p.~527]{LegRubin1987CM}.
The prime restriction in the subsequent assertion of that
theorem concerns primary triviality.  The first assertion
has no such restriction.  The original outline on page 529
explicitly states that the proof establishes finiteness
at every prime, including primes dividing the number of
CM units.  Thus the application includes two and every
odd prime, with no ordinary-reduction or prime-to-conductor
assumption added or omitted.

Apply (A.160.12) a second time, with
\(k_0=K_0,L=F,B=C,e=2\).  The two target \(\Sha\)
groups are now finite by (A.160.15) and Rubin.  The sum
of their two finite restriction images contains
\(2\Sha(C/F)\), by (A.160.12).  Hence this image is
finite.  Its kernel under multiplication by two is
\(\Sha(C/F)[2]\), already finite by the fixed-level
Selmer argument above.  Therefore \(\Sha(C/F)\) is
finite.  Isomorphism (A.160.9) gives finiteness of
\(\Sha(E_b/F)\).  Finally the kernel of restriction
\(\Sha(E_b/K)\to\Sha(E_b/F)\) lies in
\(H^1(\operatorname{Gal}(F/K),E_b(F))\), a finite
group.  This proves full finiteness of \(\Sha(E_b/K)\).
It retains A.154's exact two-primary subgroup and proves
that every odd-primary subgroup is finite and that only
finitely many are nonzero.

\emph{The smaller Weil restriction and its actual cohomology maps.}
Take \(A=\operatorname{Res}_{K/\Q}E_b\), not the auxiliary
restriction from \(F\).  For every \(\Q\)-algebra \(R\)
its defining functor is
\[
                     A(R)=E_b(K\otimes_\Q R).
\tag{A.160.16}
\]
Thus \(A(\Q)=E_b(K)\), and
\(A(\Q_p)=\prod_{w\mid p}E_b(K_w)\).
The corresponding cohomology isomorphism can be specified
on cocycles.  Write \(G=G_\Q,H_0=G_K\), and use the
coinduced module of functions \(f:G\to E_b(\overline\Q)\)
with \(f(hg)=hf(g)\), acted on by right translation.
Evaluation at the identity sends a \(G\)-cocycle \(z\)
to the \(H_0\)-cocycle \(h\mapsto z(h)(1)\).
For its inverse choose representatives \(r\) of
\(H_0\backslash G\), including one, and write
\(rg=h_{r,g}r_g\).  A cocycle \(a\) of \(H_0\)
gives the cocycle with values
\[
                            z(g)(r)=a(h_{r,g}),
\tag{A.160.17}
\]
extended to the other arguments by \(H_0\)-equivariance.
The identity
\(h_{r,g_1g_2}=h_{r,g_1}h_{r_{g_1},g_2}\)
verifies the cocycle equation.  These maps induce inverse
cohomology maps and are independent of representative
choices up to the corresponding coboundary.
On restriction to a decomposition group, the same formula
decomposes over its double cosets, namely the places of
\(K\) above that place of \(\Q\).  Hence it identifies
the localization map with the product of the localization
maps over \(K\), including real places.  Taking their
kernels proves exactly the \(\Sha\) identification in
(A.160.2).  This is the Shapiro map used here, with its
local comparison specified.

Adjunction for (A.160.16) sends the actual positive Prym
map to \(g:\overline D\to A\).  On \(\overline D_K\)
that Prym map is the displayed morphism
\[
 q_+(T,W)=\left((2\cdot191)^2\frac{aW^2}{T},
          (2\cdot191)^3\frac{aW(T+1/t)}{T}\right).
\tag{A.160.18}
\]
Over \(\overline\Q\), the map \(g\) is the pair
\((q_+,q_+^\sigma)\) into \(E_b\times E_b^\sigma\).
It is nonconstant because its first projection has degree
two.  A.150 and A.152 computed the reduced fibres
\(q_+^{-1}(O)=P_0+P_\infty\) and
\(q_+^{-1}(T_0)=P_1+P_{1/2}\).  Its conjugate has the
same fibres.  Intersecting these two explicit fibre
equations proves (A.160.3), including reducedness.

\emph{All finite Kummer levels and every physical finite place.}
Let \(s\) be a full Selmer section of \(X\) as in the
statement.  Compactification and \(g\) pull back each
multiplication torsor \([n]:A\to A\), giving compatible
classes \(\eta_n\in H^1(\Q,A[n])\).  At an ordinary
local point these are the finite Kummer classes of its
image under \(g\).  At a cusp, compactification kills
the horizontal inertia and gives the proper boundary
point, with the same Kummer assertion.  At a real place
the assertion is about the chosen real component and its
finite Kummer class.  Therefore every \(\eta_n\) belongs
to \(\operatorname{Sel}^n(A/\Q)\).

The finite multiplication sequences are
\[
 0\longrightarrow A(\Q)/nA(\Q)
  \xrightarrow{\kappa_n}\operatorname{Sel}^n(A/\Q)
  \longrightarrow\Sha(A/\Q)[n]\longrightarrow0.
\tag{A.160.19}
\]
For \(n\mid m\), the map on the right is multiplication
by \(m/n\).  Since \(\Sha(A/\Q)\) is finite, choose
\(m/n\) divisible by its exponent.  Compatibility then
makes the image of \(\eta_n\) in \(\Sha[n]\) zero.
For all even \(n\), the left groups in (A.160.19) are
the same group \(A(\Q)=\{O_A,T_A\}\), with identity
transition maps.  Their injections identify a unique
common point \(Q_A\) with
\[
                         \eta_n=\kappa_n(Q_A)
                                    \quad\text{for every }n.
\tag{A.160.20}
\]
It is enough to use the cofinal sequence of factorial
moduli; the remaining equalities follow by reduction.
This also identifies the continuous full Tate-module
class.  Indeed finite zero-cochains giving a coboundary
form finite sets.  Solutions at an arbitrarily large
factorial level solve every preceding level, so compactness
of the product of these finite sets gives a compatible
zero-cochain.  No infinitely ramified continuous-cohomology
kernel is suppressed.

For every nonarchimedean local field \(L\), the full
Kummer map on \(E_b(L)\) is injective.  Here is the
needed separation proof without a reduction hypothesis.
Choose a sufficiently small formal subgroup \(U_L\)
on which the formal logarithm identifies the group law
with an additive lattice in \(L\).  It has finite index
in the compact group \(E_b(L)\).  If \(d_L\) kills
that finite quotient, then
\(d_LE_b(L)\subset U_L\).  A point in every
\(nE_b(L)\) consequently lies in every
\(p^rU_L\), by taking \(n=d_Lp^r\); their intersection
is zero in the logarithmic lattice.  The finite Kummer
sequence identifies these subgroups as the respective
kernels.  This proves injectivity.  The product formula
in (A.160.16) proves the same assertion for \(A(\Q_p)\).
Consequently (A.160.20) forces the actual proper local
centre of \(s\) to lie in the corresponding fibre in
(A.160.3) at every finite prime.

\emph{Finite support and the cusp packet, with the neighbourhood maps.}
We use the finite-place form of Stoll's Theorem 8.2,
with the real-component correction, and the neighbourhood
argument in Stix's Corollary 5
\cite[Theorem 8.2]{LegStoll2007}
\cite[\S3]{LegStollErrata2017}
\cite[Theorem 4 and proof of Corollary 5]{BIR_Stix2015}.
Their exact input here is the fixed finite scheme
\(Z=g^{-1}(Q_A)\) in (A.160.3); it is enough that the
local centres lie in \(Z\) at a density-one set of finite
places.  We have proved this at every finite place.

For clarity the passage from an adelic point to a section
packet uses every section-neighbourhood.  Let
\((Y,s_Y)\to(X,s)\) be any geometrically connected finite
etale neighbourhood.  The map extends to the proper
normalizations.  At an ordinary local point, the
section-selected fixed sheet gives a rational local
point of the finite etale fibre.  At a cusp the selected
inertia orbit in that fibre is Galois stable; the
normalization identifies that orbit with a boundary point.
The projection of its stabilizer onto the full local
Galois group makes that boundary point rational over
the local field.  Thus the proper local centres upstairs
are supported on the finite inverse image \(Z_Y\) of
\(Z\).  Their proper genus is at least three by
Riemann--Hurwitz.

The compactified section makes this adelic tuple survive
finite descent: pulling a finite torsor back along the
section gives a global torsor class, and twisting by it
makes all the local fixed sheets rational.  This is the
actual torsor comparison of Harari--Stix, Theorem 2.1
\cite[Theorem 2.1 and unnumbered remark (1) following Theorem 2.1]{AgentHarariStix}.
It applies to arbitrary finite group schemes, and hence
to the abelian ones required by Stoll.  Stoll's theorem
therefore supplies a rational point in \(Z_Y\).
At each neighbourhood the set of rational choices is
finite and nonempty.  A common refinement supplies a
compatible choice at every finite collection of
neighbourhoods; compactness of their finite product gives
a compatible rational pro-point.

Here the pro-cover is the actual section-neighbourhood
cover \(X_s\) with arithmetic fundamental group
\(S=s(G_\Q)\); its proper normalization limit is
\(\overline X_s\).  The compatible rational points
just constructed belong to \(\overline X_s(\Q)\).
The exact neighbourhood criterion in Stix's
\S\S2.1.2--2.1.3 identifies their images in
\(\overline X(\Q)\) with the underlying points or
cusps of \(s\) \cite[\S\S2.1.2--2.1.3]{BIR_Stix2015}.
One can see the group comparison directly by prolonging
the rational pro-point through finite geometric covers.
The choices of prolongation have finite nonempty fibres;
compactness gives compatible choices.  Its decomposition
image in \(\pi_1(X_s)=S\) projects onto \(G_\Q\),
since the residue of the pro-point is \(\Q\).
The projection \(S\to G_\Q\) is an isomorphism, so
that decomposition image equals \(S\).  Under the map
to \(\pi_1(X)\) it is contained in the point
decomposition group in the interior, or the cusp
decomposition group at the boundary.  This proves the
point-section assertion in the first case.  In the second,
\(S\) meets geometric inertia trivially and is a
splitting in the rational cusp decomposition group.
Thus the packet conclusion follows from the rational
pro-point on \(X_s\), with all prolongation choices
made compatibly.  No individual rational point on a
finite neighbourhood was asserted to be fixed by that
neighbourhood's chosen section.  The possible underlying
points are precisely those in (A.160.3).

For the stated number-field classification, take any
hyperbolic open \(X/K\) in \(\overline D_K\), and use
the actual composite
\(\pi_1(X/K)\to\pi_1(\overline D_K/K)
\xrightarrow{(q_+)_*}\pi_1(E_b/K)\).
Its multiplication classes lie in
\(\operatorname{Sel}^n(E_b/K)\), with their actual local
Kummer images in \(E_b(K_w)\).  The proved finiteness
of \(\Sha(E_b/K)\) and \(E_b(K)=\{O,T_0\}\)
give (A.160.20) with \(K,E_b\) in place of \(\Q,A\).
The preceding full local Kummer injection applies at
every finite place \(w\) of \(K\), and the explicit
fibres of \(q_+\) are those in (A.160.3).  Stoll's
theorem and Stix's neighbourhood criterion have their
stated number-field scope.  Applying the same compact
finite-fibre argument to the \(K\)-section-neighbourhoods
therefore identifies an actual \(K\)-point section or
\(K\)-cusp packet at one of these four points.  This
proves the number-field assertion with no auxiliary
restriction on its real places.

For a full tripod section satisfying (A.160.4), pull
back the actual degree-eight torsor
\(W^8=t_U^2(1-t_U)/(64a^2)\).  Its class is zero,
so it has a Galois-fixed sheet.  Geometric conjugation
puts the section in \(E_D\), giving a full lift to
\(D^\circ\).  The local fixed-sheet and cusp-orbit
arguments just given prove that this lift is Selmer,
including its real component.  It is cuspidal by the
preceding argument.  The proper finite cover sends its
cusp decomposition group to the decomposition group of
its image cusp.  In the present coordinates that map is
\[
        P_0,P_1\longmapsto0,\qquad
        P_{1/2}\longmapsto1,\qquad P_\infty\longmapsto\infty.
\tag{A.160.21}
\]
The pushed-forward section is consequently cuspidal.

\emph{The all-coefficient metabelian quotient.}
The normality proof in A.157 uses the actual deck
transformations and proper boundary, and applies to
the full abelianization in (A.160.5).  Its integral
Schreier calculation, before any pro-two projection,
kills the primitive subgroup generated by
\(z=v^8\), \(\sum_{j\ {\rm even}}u_j\), and
\(\sum_{j\ {\rm odd}}u_j\).  Completing that free
integral rank-six quotient at all primes gives exactly
(A.160.6).  The element \(v\) has order eight and
acts by multiplication by \(V\), so the displayed
maps provide both the semidirect product and its inverse.
As \(\Pi/H\) is cyclic,
\(\Pi''\subset[H,H]\subset N_{\rm all}\), giving
the arithmetic map from the full metabelian quotient.

A section of this quotient in class (A.160.4) has the
same fixed sheet, hence lies in \(E_D/N_{\rm all}\)
after geometric conjugation.  Compactification and
\(q_+\), or equivalently \(g\), factor through its
proper geometric abelian kernel at every finite
multiplication level.  A local comparison conjugator
lifts to the full geometric group; the preimage of
\(E_D/N_{\rm all}\) is \(E_D\).  Thus each actual
local point or cusp section used in the comparison
has the corresponding full lift to the selected cover.
The finite Kummer classes are precisely the local
point classes used in (A.160.19).  Equations
(A.160.20) and (A.160.3), and local full-Kummer
injectivity, force its lifted centre into the boundary
at every finite place.  An ordinary tripod point has
only ordinary points in its finite etale fibre; it
therefore cannot give this localization.  This proves
the quotient assertion with its stated scope.

The literal biquadratic identities, scaling maps,
quadratic trace matrices, and adjunction fibre checks
are retained in
\path{code/harmonic_prym_cm_descent_certificate.py}
and its output.  The two new primary source scans and
their exact theorem passages are recorded in
\path{notes/harmonic_prym_cm_references.bib}.
\end{proof}
